\section{Root unblock obligation basis}
\label{sec:innerproduct-root-unblock-obligation-basis}

\begin{definition}[InnerProduct root obligation basis]
\label{def:innerproduct-root-unblock-obligation-basis}
The $\InnerProductUp$ root obligation basis is the root-facing
$\mathsf{BHist}$ package with the following ten fields.
\begin{enumerate}
  \item The traditional target is the constructive inner-product interface:
  a vector-space carrier equipped with a scalar-valued form, conjugate
  symmetry, sesquilinearity, positive-definite zero exactness, and
  orthogonality transport.
  \item The BEDC source is $\NameCert_{\VecSpaceUp}$ together with the
  scalar rows retained from $\FieldUp$.
  \item The carrier is the displayed vector predicate
  $C_V:\mathsf{BHist}\to\mathsf{Prop}$.
  \item The classifiers are the vector endpoint classifier $\sim_V$ and the
  scalar endpoint classifier $\sim_{\mathsf{scal}}$.
  \item The form row is
  $$
  \langle-,-\rangle_{\InnerProductUp}:C_V\times C_V\to C_{\mathsf{scal}}.
  $$
  \item The stability obligations are form transport, conjugate symmetry,
  sesquilinearity in the admitted arguments, positive-definite zero
  exactness, and orthogonality transport.
  \item The ledger rows are the vector source row, scalar source row, form
  endpoint row, law rows, transport rows, orthogonality row, and downstream
  consumer rows.
  \item The core theorem list consists of
  \autoref{thm:innerproduct-root-unblock-carrier-classifier},
  \autoref{thm:innerproduct-root-unblock-form-laws}, and
  \autoref{thm:innerproduct-root-unblock-downstream-threshold}.
  \item Exactness means row exhaustion: every root consumer can read only the
  rows listed in this basis, and each listed row is sourced from the displayed
  $\mathsf{BHist}$ data.
  \item There is no standard bridge field in this basis; its status is a
  paper obligation surface for later formalization.
\end{enumerate}
\end{definition}

\begin{theorem}[InnerProduct root carrier-classifier unblock]
\label{thm:innerproduct-root-unblock-carrier-classifier}
The root obligation basis of
\autoref{def:innerproduct-root-unblock-obligation-basis} exposes exactly the
carrier $C_V$, the vector classifier $\sim_V$, the scalar endpoint classifier
$\sim_{\mathsf{scal}}$, and the scalar-valued form endpoint row
$\langle-,-\rangle_{\InnerProductUp}$. It introduces no second vector
carrier, no extra point classifier, and no host scalar equality.
\end{theorem}
\begin{proof}
Read the source field of
\autoref{def:innerproduct-root-unblock-obligation-basis}. It names the
embedded $\NameCert_{\VecSpaceUp}$ and the scalar rows retained from
$\FieldUp$, which are the same source data displayed in
\autoref{def:innerproduct-bhist-source}.

Apply \autoref{thm:innerproduct-bhist-carrier-row}. This identifies the
public carrier as the displayed predicate $C_V$ together with the
scalar-valued form row, and it rules out any additional point carrier.

Apply \autoref{thm:innerproduct-classifier-transport-row}. Its hypotheses and
conclusion use exactly $\sim_V$ on vector endpoints and
$\sim_{\mathsf{scal}}$ on scalar endpoints. Hence the basis exposes precisely
the carrier, classifier, scalar endpoint classifier, and form endpoint rows
listed in the statement.
\end{proof}

\begin{theorem}[InnerProduct root form-law unblock]
\label{thm:innerproduct-root-unblock-form-laws}
For the root obligation basis of
\autoref{def:innerproduct-root-unblock-obligation-basis}, conjugate symmetry,
sesquilinearity, positive-definite zero exactness, and orthogonality are all
sourced from displayed $\mathsf{BHist}$ rows of the $\InnerProductUp$
certificate. None of these laws is supplied by a host Hilbert space, quotient
field, or ambient topology.
\end{theorem}
\begin{proof}
Use \autoref{thm:innerproduct-symmetry-positivity-row} for conjugate symmetry
and the positivity side of the diagonal law. These are scalar-history
comparisons over the retained scalar classifier.

Use \autoref{thm:innerproduct-vecspace-linearity-row} for the
sesquilinearity rows. Their vector inputs come from the embedded
$\VecSpaceUp$ operations, and their scalar outputs remain in the retained
scalar endpoint source.

Apply \autoref{thm:innerproduct-positive-definite-obligation}. It supplies
the exact zero clause saying that the self-form endpoint is scalar-zero
exactly at the displayed vector-zero classifier class.

Finally unpack \autoref{def:innerproduct-orthogonality-row}. Orthogonality is
the scalar-zero classification of the displayed form endpoint, so its
transport and zero behavior are inherited from the same form and classifier
rows. Thus every form law named by the basis is a displayed
$\mathsf{BHist}$ row.
\end{proof}

\begin{theorem}[InnerProduct root downstream threshold unblock]
\label{thm:innerproduct-root-unblock-downstream-threshold}
The downstream certificates $\HilbertUp$, $\NormUp$, $\FourierUp$,
$\RiemannianMetricUp$, and $\RootSystemUp$ may consume from
\autoref{def:innerproduct-root-unblock-obligation-basis} only the carrier,
classifier, scalar endpoint, form-law, orthogonality, and ledger rows listed
there. They receive no Banach completeness, Cauchy-limit selector, ambient
topology, or host Hilbert-space row from $\InnerProductUp$.
\end{theorem}
\begin{proof}
Start with \autoref{def:innerproduct-root-threshold-source}. It packages the
carrier, classifier, scalar endpoint source, form row, form-law rows,
orthogonality row, and ledger rows as the complete inner-product threshold
surface.

Apply \autoref{thm:innerproduct-root-threshold-orthogonality-consumption}.
This separates the inner-product geometric rows consumed by $\HilbertUp$
from the Banach metric, completeness, and Cauchy-limit rows supplied by the
separate Banach dependency.

Apply \autoref{thm:innerproduct-root-threshold-downstream-coverage}. It gives
the same threshold coverage for $\HilbertUp$, $\FourierUp$,
$\RiemannianMetricUp$, and $\RootSystemUp$, and it excludes Banach
completeness, Cauchy-limit selection, ambient topology, and host Hilbert-space
data from the $\InnerProductUp$ source.

The $\NormUp$ consumer reads the same inner-product threshold through the norm
and metric seed rows of \autoref{def:innerproduct-norm-metric-seed-source}.
Those rows are generated from the form endpoint, positive-definite
exactness, classifier transport, and ledger fields already listed in the
basis. Therefore all five consumers are bounded by the same root obligation
basis, with no additional downstream row projected from $\InnerProductUp$.
\end{proof}
